% ============================================================
% TEMA E APARÊNCIA
% ============================================================
\usetheme{CambridgeUS}
\usecolortheme{default}
\setbeamertemplate{headline}{}
\setbeamertemplate{navigation symbols}{}
\setbeamertemplate{bibliography item}{}
\setbeamertemplate{caption}[numbered]
\setbeamertemplate{itemize item}[square]
\setbeamertemplate{itemize subitem}[square]
\setbeamertemplate{itemize subsubitem}[square]
\setbeamercolor{title}{use=author in head/foot, bg=author in head/foot.bg, fg=white}
\setbeamercolor{frametitle}{use=author in head/foot, bg=author in head/foot.bg, fg=white}
\setbeamercolor{normal text}{fg=black}
\setbeamercolor{structure}{fg=black}
\setbeamertemplate{blocks}[rounded][shadow=false]
\setbeamercolor{block title}{use=author in head/foot,fg=black,bg=author in head/foot.bg!20}
\setbeamercolor{block body}{use=author in head/foot,fg=black,bg=white}
\setbeamerfont{block title}{series=\bfseries}
\setbeamerfont{frametitle}{series=\bfseries}
\setbeamercolor{itemize item}{fg=black}
\setbeamercolor{itemize subitem}{fg=black}
\setbeamercolor{itemize subsubitem}{fg=black}
\setbeamercolor{enumerate item}{fg=black}
\setbeamercolor{enumerate subitem}{fg=black}
\setbeamercolor{enumerate subsubitem}{fg=black}
\setbeamercolor{section number projected}{bg=black, fg=white}
\setbeamercolor{subsection number projected}{bg=black, fg=white}

% ============================================================
% PACOTES
% ============================================================

% Codificação e idioma
\usepackage[utf8]{inputenc}
\usepackage[T1]{fontenc}
\usepackage[brazil]{babel}
\usepackage{multirow}

% Matemática
\usepackage{amsmath}
\usepackage{amssymb}

% Gráficos e tabelas
\usepackage{graphicx}
\usepackage{booktabs}
\usepackage{threeparttable}

% Referências no padrão ABNT
\usepackage[alf,bibjustif,recuo=0cm,abnt-etal-cite=4,abnt-etal-text=it]{abntex2cite}

% ============================================================
% COMANDOS CUSTOMIZADOS
% ============================================================

\newcounter{citeabntcount}

% Auxiliar interno — não usar diretamente
\newcommand{\citeabntitem}[1]{%
    \stepcounter{citeabntcount}%
    \ifnum\value{citeabntcount}>1\relax; \fi
    \citeauthoronline{#1}, \citeyear{#1}%
}

% Citação ABNT: \citeabnt[p. X]{chave1,chave2}
\newcommand{\citeabnt}[2][]{%
    \begingroup
    \setcounter{citeabntcount}{0}%
    (\forcsvlist{\citeabntitem}{#2}\ifstrempty{#1}{}{, p.~#1})%
    \endgroup
}
